\documentclass[../LogosReference.tex]{subfiles}
\begin{document}

% ============================================================================
% Section: Introduction
% ============================================================================

\section{Introduction}\label{sec:introduction}

This reference manual provides the formal specification for the Logos proof system, semantics, metalogic, and library of theorems.
The semantics is organized around the semantic frames that provides the primitive structure used to interpret the logical operators included in the Logos.
By successively extending the primitive resources included in a frame, it becomes possible to provide semantic clauses for a new range of operators, thereby extending the expressive power of the Logos for systematic reasoning. 

The Logos is interpreted over a hyperintensional \textit{state semantics} which comprises the \textit{Constitutive Foundation}.
By contrast to intensional semantic theories which take worlds to be primitive, the state semantics takes states and their parthood relations to be primitive to provide fine-grained semantic contents by which to recursively interpret the primitive expressions of the language.
The \textit{Dynamical Foundation} extends the state semantics to include a \textit{task relation} parameterized by durations belonging to a \textit{temporal order} in order to govern the lawful transitions by which a system and its parts evolve.
These ingredients are then used to construct \textit{world sates} which represent total instantaneous states of the system under study, as well as \textit{world histories} which model the possible evolutions of that system by mapping durations to world states.
For instance, in modeling a game of chess, the world states correspond to the legal complete configurations of the chessboard, and the world histories correspond to the possible games of chess.

Given a world history, each duration may be referred to as a \textit{time} in that history.
Together, a world history and time suffice to recursively assign truth-values to all sentences of the foundational layers of the Logos.
In addition to the standard Boolean and first-order operators, the foundational semantic layers for the Logos support modal, temporal, counterfactual, constitutive, and causal operators for complex reasoning including planning and evaluating past and future actions.
Modular plugins are then provided to extend the expressive power of the foundational layers of the Logos to include sets of related logical operators.
In addition to extending the semantic primitives included in a frame, plugins may contribute new contextual parameters at which sentences are evaluated.
For instance, given the addition of a \textit{credence function} which assigns likelihoods to propositions, a \textit{belief parameter} is required along with a world history and time to provide a semantics for the belief operator.
Given a world history, time, and belief parameter, a proposition is believed just in case the credence in that proposition is greater than the belief parameter included in the point of evaluation.


% ------------------------------------------------------------
% Extension Dependencies
% ------------------------------------------------------------

\subsection{Extension Dependencies}\label{sec:extension-dependencies}

This section presents the layer extensions by which the Logos is constructed where \Cref{fig:extension-dependencies} shows the dependency structure among extensions where the first two layers provide an essential foundation followed by a modular plugin system for language extensions.

\begin{figure}[htbp]
\centering
\usetikzlibrary{fit,backgrounds,positioning,calc,arrows.meta}
\begin{tikzpicture}[
  widebox/.style={rectangle, draw, text centered, rounded corners=4pt, minimum width=10cm, minimum height=1.2cm},
  box/.style={rectangle, draw, text centered, rounded corners=4pt, minimum width=4.5cm, minimum height=1.2cm},
  smallbox/.style={rectangle, draw, text centered, rounded corners=3pt, minimum width=2.8cm, minimum height=1.6cm, font=\small},
  arrow/.style={->, thick, >=stealth},
  doublearrow/.style={<->, thick, >=stealth}
]

% Constitutive Foundation (full-width)
\node[widebox, minimum height=1.8cm] (foundation) {
  \begin{tabular}{c}
    \textbf{Constitutive Foundation} \\
    {\small predicates ($F$, $G$),~ functions ($f$, $g$),~ variables ($x$, $y$)} \\
    {\small boolean ($\neg$, $\land$, $\lor$, $\bot$, $\top$),~ lambda ($\lambda$),~ identity ($=$)} \\
    {\small propositional identity ($\equiv$),~ grounding ($\leq$),~ essence ($\sqsubseteq$),~ reduction ($\Rightarrow$)}
  \end{tabular}
};

% Explanatory Extension (full-width)
\node[widebox, minimum height=1.8cm, below=0.85cm of foundation] (explanatory) {
  \begin{tabular}{c}
    \textbf{Explanatory Extension} \\
    {\small modal ($\Box$, $\Diamond$),~ temporal ($H$, $G$, $P$, $F$, $S$, $U$),~ store/recall ($\downarrow$, $\uparrow$)} \\
    {\small counterfactual ($\boxright$, $\diamondright$),~ causal ($\circleright$, $\dotcircleright$)} \\
    {\small quantifiers ($\forall$, $\exists$),~ actuality ($\mathsf{Act}$)}
  \end{tabular}
};

% Middle Extensions (two rows: 3 on top, 2 centered below)
% Top row center: Normative
\node[smallbox, below=2.0cm of explanatory] (normative) {
  \begin{tabular}{c}
    \textbf{Normative} \\
    {\footnotesize obligation, permission,} \\
    {\footnotesize preference, explanation} \\
    {\small ($\mathsf{O}$, $\mathsf{P}$, $\prec$, $\mapsto$)}
  \end{tabular}
};

% Top row left: Epistemic
\node[smallbox, left=0.3cm of normative] (epistemic) {
  \begin{tabular}{c}
    \textbf{Epistemic} \\
    {\footnotesize belief, probability,} \\
    {\footnotesize might/must} \\
    {\footnotesize indicative} \\
    {\small ($\mathsf{B}$, $\mathsf{Pr}$, $\mathsf{Mi}$, $\mathsf{Mu}$, $\hookrightarrow$)}
  \end{tabular}
};

% Top row right: Abilities
\node[smallbox, right=0.3cm of normative] (abilities) {
  \begin{tabular}{c}
    \textbf{Abilities} \\
    {\footnotesize specific, generic,} \\
    {\footnotesize agency} \\
    {\small ($\mathsf{Can}$, $\mathsf{Stit}$)}
  \end{tabular}
};

% Bottom row: Choice and Spatial centered below
\node[smallbox, below=0.5cm of normative, xshift=-1.6cm] (choice) {
  \begin{tabular}{c}
    \textbf{Choice} \\
    {\footnotesize free permission,} \\
    {\footnotesize alternatives} \\
    {\small ($\mathsf{FP}$, $\mathsf{FF}$, $\mathsf{Ch}$)}
  \end{tabular}
};

\node[smallbox, below=0.5cm of normative, xshift=1.6cm] (spatial) {
  \begin{tabular}{c}
    \textbf{Spatial} \\
    {\footnotesize location,} \\
    {\footnotesize relations} \\
    {\small ($@$, $\mathsf{Btw}$, \ldots)}
  \end{tabular}
};

% Ellipsis node for extensibility (second row)
\node[draw=none, right=0.3cm of spatial, font=\Large] (rightellipsis) {$\cdots$};

% Invisible spacer nodes to extend grey box
\node[draw=none, above=0.7cm of normative, minimum width=1pt, minimum height=1pt] (topspacer) {};
\node[draw=none, below=0.0cm of choice, minimum width=1pt, minimum height=1pt] (bottomspacer) {};

% Grey background box around middle extensions
\begin{pgfonlayer}{background}
  \node[fill=gray!15, draw=gray!40, rounded corners=6pt,
        fit=(epistemic)(normative)(abilities)(rightellipsis)(choice)(spatial)(topspacer)(bottomspacer),
        inner sep=12pt,
        label={[font=\itshape, anchor=north, yshift=-8pt]north:Middle extensions are modular---combine any subset}] (middlebox) {};
\end{pgfonlayer}

% Bottom Layer: Agential, Social, and Reflection (inside third grey box)
% Position them below the middle extensions grey box, centered
\node[smallbox, below=2.0cm of middlebox] (social) {
  \begin{tabular}{c}
    \textbf{Social} \\
    {\footnotesize common ground,} \\
    {\footnotesize trust relations}
  \end{tabular}
};

\node[smallbox, left=0.3cm of social] (agential) {
  \begin{tabular}{c}
    \textbf{Agential} \\
    {\footnotesize inherited agent-indexed} \\
    {\footnotesize operators}
  \end{tabular}
};

\node[smallbox, right=0.3cm of social] (reflection) {
  \begin{tabular}{c}
    \textbf{Reflection} \\
    {\footnotesize metacognition,} \\
    {\footnotesize perspectival operators}
  \end{tabular}
};

% Invisible spacer node to extend grey box upward for label
\node[draw=none, above=0.5cm of social, minimum width=1pt, minimum height=1pt] (bottomtopspacer) {};

% Grey background box around bottom extensions
\begin{pgfonlayer}{background}
  \node[fill=gray!15, draw=gray!40, rounded corners=6pt,
        fit=(agential)(social)(reflection)(bottomtopspacer),
        inner sep=12pt,
        label={[font=\itshape, anchor=north, yshift=-8pt]north:Agent-dependent extensions}] (bottombox) {};
\end{pgfonlayer}

% Simplified arrows (no labels)
% Foundation -> Explanatory
\draw[arrow] (foundation) -- (explanatory);

% Explanatory -> Middle box (arrow from center of explanatory, straight down)
\draw[arrow] (explanatory.south) -- (explanatory.south |- middlebox.north);

% Middle box -> Bottom box (single centered arrow)
\draw[arrow] (middlebox.south) -- (middlebox.south |- bottombox.north);

% Double arrow between Agential and Reflection
% \draw[doublearrow] (agential.east) -- (reflection.west);

\end{tikzpicture}
\caption{Extension dependency structure in the Logos system.}
\label{fig:extension-dependencies}
\end{figure}

Whereas the Constitutive Foundation and Explanatory Extension form the required base, the middle extensions (Epistemic, Abilities, Normative, Choice, Spatial, and others) are modular plugins that can be combined in any subset.
The Agential Extension, Social Extension, and Reflection Extension all inherit from the layers above and are free to interact with each other.
The Agential Extension provides third-person multi-agent reasoning, the Social Extension provides the construction of the common ground and trust relations, and the Reflection Extension enables first-person metacognition.

% ------------------------------------------------------------
% Layer Descriptions
% ------------------------------------------------------------

\subsection{Layer Descriptions}\label{sec:layer-descriptions}

\begin{enumerate}
  \item \textbf{Constitutive Foundation}:
        Provides the hyperintensional base with predicates ($F$, $G$), functions ($f$, $g$), variables ($x$, $y$), booleans ($\neg$, $\land$, $\lor$), extremal ($\bot$, $\top$), lambda abstraction ($\lambda$), identity ($=$), and constitutive operators (propositional identity $\equiv$, grounding $\leq$, essence $\sqsubseteq$, reduction $\Rightarrow$) over a complete lattice with bilateral propositions (verifier/falsifier pairs).
        This enables fine-grained distinctions between propositions with identical truth-values and modal profiles but different verification conditions, essential for exact specification of what makes a claim true.

  \item \textbf{Dynamical Foundation}:
        Adds quantifiers ($\forall$, $\exists$), modals ($\Box$, $\Diamond$) for necessity and possibility, tense ($H$, $G$, $P$, $F$, $\since$, $\until$) for reasoning across time, store/recall operators ($\downarrow$, $\uparrow$) for temporal reference, actuality ($\mathsf{Act}$) for distinguishing actual from possible, counterfactual conditionals ($\boxright$, $\diamondright$) for hypothetical reasoning, and causal operators ($\circleright$, $\dotcircleright$) for explanatory links.
        Together these resources enable reasoning about action sequences, future contingency, what would have happened under different circumstances, and causal relationships.

  \item \textbf{Epistemic Extension}:
        Introduces epistemic modals ($\mathsf{Mi}$, $\mathsf{Mu}$), belief ($\mathsf{B}$), probability ($\mathsf{Pr}$), and indicative conditional operators ($\hookrightarrow$) for reasoning under uncertainty.
        By extending the semantic framework with credence functions to assign weights to state transitions, agents can track how evidence updates beliefs across time and distinguish what an agent believes from what is actually the case.

  \item \textbf{Abilities Extension}:
        Provides ability ($\mathsf{Can}$) and STIT ($\mathsf{stit}$) operators for reasoning about capacities.
        Ability modals express what an agent can bring about through their capacities, with stronger truth conditions than circumstantial possibility since ability entails possibility but not \textit{vice versa}.
        This enables reasoning about what agents are capable of accomplishing given their intrinsic capacities and circumstances.

  \item \textbf{Normative Extension}:
        Adds obligation ($\mathsf{O}$), permission ($\mathsf{P}$), and a preference operator ($\prec$) with value orderings over states for reasoning about what ought to be done, what is permitted, and how to rank alternatives.
        This enables ethical reasoning about alternatives.

  \item \textbf{Choice Extension}:
        Introduces free choice permission ($\mathsf{FP}$) and best choice ($\mathsf{Ch}$) operators for reasoning about permission over alternatives, enabling reasoning about which alternatives are permitted and how permissions combine.

  \item \textbf{Spatial Extension}:
        Provides location-dependent operators and spatial relations for reasoning about where things are and they are oriented relative to an orientation basis.
        This supports navigation, containment reasoning, and location-sensitive planning.

  \item \textbf{Agential Extension}:
        Indexes epistemic and normative operators to specific agents ($\mathsf{B}_a$, $\prec$, $\mathsf{O}_a$, \ldots), enabling reasoning about what different agents believe, prefer, and are obligated to do.
        This transforms single-perspective reasoning into multi-agent coordination where agents can model each other's epistemic and normative states.

  \item \textbf{Social Extension}:
        Provides resources for common ground construction and group-level reasoning, building upon the Agential Extension to enable reasoning about shared beliefs, build trust in other agents, and coordinate on common goals.

  \item \textbf{Reflection Extension}:
        Adds first-person metacognitive operators ($\mathsf{I}$) for self-directed reasoning about one's own beliefs, abilities, and progress towards goals.
        This enables an agent to reason about what it is in a good position to believe given evaluation of its own limitations and track relative to other agents. Essential for AI systems that must understand their own capabilities and benefit from the capacity to learn by the testimony of others.
\end{enumerate}

% ------------------------------------------------------------
% Document Organization
% ------------------------------------------------------------

\subsection{Document Organization}\label{sec:document-organization}

This reference manual is organized as follows:

\begin{description}
  \item[\Cref{sec:constitutive-foundation}] presents the Constitutive Foundation, including the mereological state space, verification and falsification clauses, and bilateral propositions.

  \item[\Cref{sec:core-syntax}] introduces the syntactic primitives of the Explanatory Extension, including modal, temporal, and counterfactual operators.

  \item[\Cref{sec:core-semantics}] provides the semantic framework for the Explanatory Extension, including core frames, world-histories, and truth conditions.

  \item[\Cref{sec:core-axioms}] presents the axiom system for counterfactual logic.\\

  \ldots\\

  \item[\Cref{sec:reflection}] specifies the Reflection Extension for first-person metacognitive reasoning, including truth conditions, derived operators, and temporal interactions.
\end{description}


% ------------------------------------------------------------
% Lean Implementation
% ------------------------------------------------------------

\subsection{Lean Implementation}\label{sec:lean-implementation}

The Logos system is implemented in Lean 4 with Mathlib. Each definition in this reference has a corresponding Lean implementation, referenced using the \verb|\leansrc{}{}| command.

For example, \leansrc{Logos.Foundation.Frame}{ConstitutiveFrame} refers to the \texttt{ConstitutiveFrame} definition in the \texttt{Logos.Foundation.Frame} module.
The implementation can be found in the \proofchecker{} repository in the \texttt{Theories/Logos/} directory.

\end{document}
